\documentclass[11pt]{article}
\usepackage[a4paper,margin=1in]{geometry}
\usepackage{fontspec}
\usepackage[boldfont=false]{xeCJK}
\setCJKmainfont{FandolSong-Regular.otf}
\usepackage{amsmath}
\usepackage{amssymb}
\usepackage{array}
\usepackage{booktabs}
\usepackage{graphicx}
\usepackage{adjustbox}
\usepackage{enumitem}
\setlist[itemize]{topsep=2pt,partopsep=0pt,parsep=0pt,itemsep=2pt,leftmargin=1.4em}
\setlist[enumerate]{topsep=2pt,partopsep=0pt,parsep=0pt,itemsep=2pt,leftmargin=1.6em}
\usepackage{parskip}
\usepackage{fvextra}
\fvset{breaklines=true,breakanywhere=true,fontsize=\small}
\setlength{\parindent}{0pt}

\begin{document}

\begin{center}
  {\LARGE\bfseries Data representation}\\[0.35em]
  {\large Igcse Computer Science}
\end{center}
\vspace{0.6em}

\section*{Why computers use binary}

A computer can only work with two states: on and off. You write these as 1 and 0. A system that uses only two digits 数字 is called \textbf{binary} 二进制 (base 2).

Every kind of data 数据 — numbers, text, sound and images — must be changed into binary before a computer can use it. The computer processes this binary using \textbf{logic gates} 逻辑门, and stores it in \textbf{registers} 寄存器 (small, fast stores inside the processor 处理器).

\section*{Number systems}

A \textbf{number system} 数制 is a way of writing numbers using a fixed set of digits. You need three of them.

\begin{center}
\adjustbox{max width=\linewidth}{%
\begin{tabular}{lll}
\toprule
\textbf{System} & \textbf{Base} & \textbf{Digits used} \\
\midrule
Denary & 10 & 0–9 \\
Binary & 2 & 0 and 1 \\
Hexadecimal & 16 & 0–9 then A–F \\
\bottomrule
\end{tabular}}
\end{center}

\begin{itemize}
\item \textbf{denary} 十进制 is the normal counting system (also called decimal).

\item \textbf{binary} uses only 0 and 1.

\item \textbf{hexadecimal} 十六进制 (hex) uses sixteen digits: 0–9, then A, B, C, D, E, F stand for 10, 11, 12, 13, 14, 15.

\end{itemize}

The \textbf{base} 基数 tells you how many different digits a system uses.

\subsection*{Place value}

Each column in a number has a \textbf{place value} 位值. In binary the place values double from right to left. For an 8-bit number they are:

\begin{Verbatim}
128  64  32  16  8  4  2  1
\end{Verbatim}

One \textbf{bit} 位 is a single 0 or 1. Eight bits make one \textbf{byte} 字节. Four bits (half a byte) is a \textbf{nibble} 半字节.

\subsection*{Converting between number systems}

\textbf{Denary → binary.} Write the place values. Put a 1 under each value you need so they add up to your number; put 0 under the rest.

Example: change denary 150 to binary. $150 = 128 + 16 + 4 + 2$.

\begin{Verbatim}
128 64 32 16 8 4 2 1
  1  0  0  1 0 1 1 0
\end{Verbatim}

So $150$ = \texttt{10010110}.

\textbf{Binary → denary.} Add the place values where there is a 1. \texttt{10010110} $= 128 + 16 + 4 + 2 = 150$.

\textbf{Hexadecimal → binary.} Change each hex digit into its own 4-bit group (a nibble).

Example: hex F08. $F = 1111$, $0 = 0000$, $8 = 1000$, so F08 = \texttt{1111 0000 1000}.

\textbf{Binary → hexadecimal.} Group the bits into nibbles of 4, starting from the right. Change each nibble to one hex digit.

\textbf{Denary → hexadecimal.} The easy way is to change to binary first, then binary to hex.

This table helps with the hex letters:

\begin{center}
\adjustbox{max width=\linewidth}{%
\begin{tabular}{lll}
\toprule
\textbf{Denary} & \textbf{Binary} & \textbf{Hex} \\
\midrule
10 & 1010 & A \\
11 & 1011 & B \\
12 & 1100 & C \\
13 & 1101 & D \\
14 & 1110 & E \\
15 & 1111 & F \\
\bottomrule
\end{tabular}}
\end{center}

Cambridge questions use binary numbers up to 16 bits long.

\subsection*{Why hexadecimal is used}

Hex is shorter than binary and easier for people to read and write. One hex digit replaces 4 binary digits, so you make fewer mistakes. The value does not change — hex is just a shorter way to show the same binary.

Computer scientists use hex for:

\begin{itemize}
\item MAC addresses and IPv6 addresses

\item colour codes in HTML (for example \texttt{\#FF0000} is red)

\item \textbf{memory addresses} 内存地址 and error codes

\item showing the contents of memory (a "memory dump")

\end{itemize}

\section*{Binary addition}

You can add two 8-bit binary numbers, column by column from the right, just like denary. The rules for one column are:

\begin{center}
\adjustbox{max width=\linewidth}{%
\begin{tabular}{llll}
\toprule
\textbf{A} & \textbf{B} & \textbf{Result bit} & \textbf{Carry} \\
\midrule
0 & 0 & 0 & 0 \\
0 & 1 & 1 & 0 \\
1 & 0 & 1 & 0 \\
1 & 1 & 0 & 1 \\
\bottomrule
\end{tabular}}
\end{center}

When a carry also comes in, $1 + 1 + 1 = 1$ with a carry of 1.

Example: add \texttt{01110110} (118) and \texttt{00110000} (48).

\begin{Verbatim}
  0 1 1 1 0 1 1 0    (118)
+ 0 0 1 1 0 0 0 0    (48)
-------------------
  1 0 1 0 0 1 1 0    (166)
\end{Verbatim}

\subsection*{Overflow}

An 8-bit register can hold denary values from 0 to 255 only. If an addition gives a result above 255, the answer needs a 9th bit. The register cannot hold this extra bit, so it is lost. This is called \textbf{overflow} 溢出 (an overflow error). It happens when a value goes outside the limit the register can store.

Example: \texttt{11001000} (200) $+$ \texttt{01001000} (72) $= 272$. In binary that is \texttt{1 00010000}, which needs 9 bits. The leading 1 will not fit in 8 bits, so the stored answer is wrong.

\section*{Logical binary shift}

A \textbf{logical binary shift} 逻辑二进制移位 moves all the bits left or right by a number of places.

\begin{itemize}
\item Bits that move off the end of the register are \textbf{lost}.

\item \textbf{Zeros} are added at the empty end.

\end{itemize}

A \textbf{left shift} multiplies the number by 2 for each place moved. A \textbf{right shift} divides it by 2 for each place; the right-most bits (the \textbf{least significant bit(s)} 最低有效位) are lost.

Example: left shift \texttt{00110101} (53) by 2 places.

\begin{Verbatim}
start:         0 0 1 1 0 1 0 1
left shift 2:  1 1 0 1 0 1 0 0
\end{Verbatim}

The result is \texttt{11010100} (212), which is $53 \times 4$. The two left-most bits were lost and two zeros came in on the right. If a 1 is pushed off the end, that information is gone for good.

\section*{Two's complement}

So far the numbers were positive. \textbf{Two's complement} 补码 lets an 8-bit register hold negative numbers too.

In two's complement, the left-most bit (the \textbf{most significant bit} 最高有效位, or MSB) has a \emph{negative} place value:

\begin{Verbatim}
-128  64  32  16  8  4  2  1
\end{Verbatim}

\begin{itemize}
\item If the MSB is 0, the number is positive.

\item If the MSB is 1, the number is negative.

\end{itemize}

\textbf{To make a positive number negative:} write the positive binary, flip every bit (0↔1), then add 1.

Example: make $-40$.

\begin{itemize}
\item $+40$ = \texttt{00101000}

\item flip the bits = \texttt{11010111}

\item add 1 = \texttt{11011000}

\end{itemize}

So $-40$ = \texttt{11011000}. Check by adding the place values: $-128 + 64 + 16 + 8 = -40$.

To read a negative two's complement number, just add the place values (the MSB counts as $-128$). The range of an 8-bit two's complement number is $-128$ to $+127$.

\section*{Representing text}

Computers store text by giving every character a number, then storing that number in binary. The set of characters a computer can use, together with their numbers, is a \textbf{character set} 字符集.

\begin{itemize}
\item \textbf{ASCII} uses 7 bits per character, so it has 128 different characters. This is enough for English letters, digits and common symbols.

\item \textbf{Unicode} uses more bits per character. It can represent far more characters — many languages, plus symbols and \textbf{emoji} 表情符号.

\end{itemize}

Because Unicode has more characters, it needs more bits per character than ASCII, so the same text takes more storage 存储.

\section*{Representing sound}

A \textbf{sound wave} 声波 is smooth and always changing. To store it, the computer measures the height of the wave at regular moments. This is called \textbf{sampling} 采样, and each measurement is a sample.

\begin{itemize}
\item \textbf{sample rate} 采样率 is the number of samples taken each second (measured in Hz).

\item \textbf{sample resolution} 采样分辨率 is the number of bits used for each sample. The height of the wave at a sample point is its \textbf{amplitude} 振幅.

\end{itemize}

A higher sample rate and a higher sample resolution give a more accurate recording, but a larger file.

\section*{Representing images}

A computer image is made of a grid of small dots called \textbf{pixels} 像素.

\begin{itemize}
\item \textbf{resolution} 分辨率 is the number of pixels in the image (for example $1920 \times 1080$).

\item \textbf{colour depth} 颜色深度 is the number of bits used to store the colour of each pixel.

\end{itemize}

A higher resolution and a higher colour depth give a better-quality image, but a larger file.

\section*{Measuring data storage}

Data storage is measured in the units below. Each unit is 1024 times the one before it (because $1024 = 2^{10}$, which fits binary).

\begin{center}
\adjustbox{max width=\linewidth}{%
\begin{tabular}{ll}
\toprule
\textbf{Unit} & \textbf{Equals} \\
\midrule
bit & a single 0 or 1 \\
nibble & 4 bits \\
byte & 8 bits \\
kibibyte (KiB) & 1024 bytes \\
mebibyte (MiB) & 1024 KiB \\
gibibyte (GiB) & 1024 MiB \\
tebibyte (TiB) & 1024 GiB \\
pebibyte (PiB) & 1024 TiB \\
exbibyte (EiB) & 1024 PiB \\
\bottomrule
\end{tabular}}
\end{center}

\section*{Calculating file size}

\textbf{Image file size} (in bits) $=$ resolution $\times$ colour depth $=$ width $\times$ height $\times$ colour depth.

Example: an image is $1024 \times 1024$ pixels with a colour depth of 2 bytes ($= 16$ bits).

\begin{itemize}
\item bits $= 1024 \times 1024 \times 16 = 16\,777\,216$ bits

\item bytes $= \div 8 = 2\,097\,152$ bytes

\item KiB $= \div 1024 = 2048$ KiB

\item MiB $= \div 1024 = 2$ MiB

\end{itemize}

\textbf{Sound file size} (in bits) $=$ sample rate $\times$ sample resolution $\times$ length in seconds.

Always divide by 1024 (not 1000) to change to KiB, MiB and so on. Give your answer in the unit the question asks for.

\section*{Compression}

\textbf{Compression} 压缩 makes a file smaller. A smaller file:

\begin{itemize}
\item uses less storage space,

\item needs less \textbf{bandwidth} 带宽 (the amount of data a connection can carry),

\item takes a shorter time to send (a shorter \textbf{transmission} 传输 time).

\end{itemize}

There are two types.

\subsection*{Lossless compression}

\textbf{Lossless} 无损 compression makes the file smaller with \textbf{no permanent loss} of data. The original file can be rebuilt exactly.

One method is \textbf{run-length encoding} 行程编码 (RLE). It replaces a run of repeated values with one copy of the value and a count of how many times it repeats. For example \texttt{WWWWWWWW} (8 whites) is stored as "8 W". This works well when data has many repeats.

\subsection*{Lossy compression}

\textbf{Lossy} 有损 compression makes the file much smaller by \textbf{permanently removing} some data. The removed data cannot be got back. For example:

\begin{itemize}
\item reducing the resolution or colour depth of an image,

\item reducing the sample rate or sample resolution of a sound.

\end{itemize}

Use lossless when you must keep every detail (text and program files). Use lossy for photos, music and video, where a small loss of quality is worth a much smaller file.


\end{document}
